\documentclass[varwidth, border = 3pt]{standalone}

\usepackage{main}

\begin{document}

\paragraph{Exemple 1.}

%%%
% La méthode demande de régler à la main la commande ''\nubdepth'',
% valant ''0'' initialement, afin d'indiquer la profondeur du trait
% le plus bas. La valeur donnée est simple, c'est le successeur du
% nombre de macros ''\nub'' utilisées (ceci laisse envisager une
% automatisation possible).
%%%
\renewcommand{\nubdepth}{6}

$ \nub{1}{
    11 + \nub{2}{
      12 + \nub{3}{
        13 + \nub{4}{
          14 + \nub{5}{15 + 16} +
          17} +
        18} +
      19} +
    20}
= 31 \times 5
= 155$

\paragraph{Exemple 2.}

%%%
% Ce second exemple implique la mise à jour de la valeur de la
% commande ''\nubid'' nécessaire en coulisse pour nommer des
% marques nécessaires au tracé (en coulisse, un compteur pourrait
% faire ce travail à notre place).
%%%
\renewcommand{\nubid}{usecase-B}
\renewcommand{\nubdepth}{4}

$ \nub{1}{
    1 + \nub{2}{
      3 + \nub{3}{
        5
          + \cdots + % Un peu de maths... Au milieu, nous avons
        95} +        % `49 + 51`, d'où `(49 - 1) / 2 = 24` ajouts
      97} +          % du nombre `2` pour aller de 1 à 49.
    99}
= 100 \times (1 + 24)
= 2500$

\end{document}
